%
\subsection*{Variables, datos y flujo}
\label{CH10-B01}
En $u(x,t)$, $u_t$ deriva con $x$ fijo y $u_x$ con $t$ fijo. Al integrar
respecto de $x$, la constante puede depender de $t$. Los modelos básicos
son
\[
 u_t=\kappa u_{xx}\quad\text{(calor)},\qquad
 u_{tt}=c^2u_{xx}\quad\text{(onda)},\qquad
 u_{xx}+u_{yy}=0\quad\text{(equilibrio)}.
\]
La difusividad $\kappa>0$ tiene unidades de longitud$^2$/tiempo y
$c>0$ es una velocidad. Calor requiere un dato inicial temporal; onda,
desplazamiento y velocidad. Los contornos restringen el estado o su
derivada en el borde espacial.

Con densidad volumétrica de energía $e$ y flujo $J$ hacia la derecha,
$e_t+J_x=0$ implica
$\frac{d}{dt}\int_a^b e\,dx=J(a)-J(b)$, sin fuentes y por unidad de
sección. La ley de Fourier es $J=-k u_x$; el flujo saliente cambia de
signo entre extremos. Temperatura prescrita y potencia prescrita son
datos distintos; para pasar de flujo a potencia se necesita la sección.

%
\subsection*{Separación y perfiles de evolución}
\label{CH10-B02}
Con extremos homogéneos, se ensaya $u=X(x)T(t)$ y se elige la base
espacial a partir de los contornos. Para $X''=-\lambda X$, cada modo
de calor tiene amplitud $T'= -\kappa\lambda T$ y cada modo de onda
$T''=-c^2\lambda T$. Los coeficientes se obtienen proyectando los datos
iniciales con la norma de la base. Para $\lambda>0$ y
$\omega=c\sqrt\lambda$, la onda usa
$T=A\cos\omega t+(B/\omega)\sin\omega t$, donde $A$ procede del
desplazamiento y $B$ de la velocidad. Si se admite un modo cero,
su amplitud es constante en calor y afín en el tiempo en onda.

En la recta, la fórmula de d'Alembert para datos $f=u(\cdot,0)$ y
$g=u_t(\cdot,0)$ es
\[
 u(x,t)=\frac{f(x-ct)+f(x+ct)}2
       +\frac1{2c}\int_{x-ct}^{x+ct}g(s)\,ds.
\]
Con datos suficientemente regulares satisface la ecuación clásica; los
perfiles a trozos requieren precisar los puntos de regularidad.
En un intervalo no se ignoran las reflexiones en los extremos.
Una cuerda recta puede tener velocidad no nula. Distínganse perfil a
tiempo fijo e historia de un punto, y compruébese si el eje representa
$u$, $cu$ u otra escala. Consultas: soluciones de \exref{M2-E0264} y
\exref{M2-E0107}.

%
\subsection*{Equilibrio y homogeneización del contorno}
\label{CH10-B03}
Con temperaturas de borde constantes, búsquese primero un perfil
estacionario $v(x)$ que satisfaga esos contornos; para calor sin fuentes,
$v''=0$. Entonces $w=u-v$ tiene contornos homogéneos y dato inicial
$u(x,0)-v(x)$. Con aislamiento en ambos extremos se admite un modo
constante y se conserva la temperatura media en el modelo homogéneo.

En un rectángulo $0<x<L$, $0<y<H$, con tres lados a cero y dato en
$y=H$, los modos de Laplace utilizan
\[
 \sin\frac{n\pi x}{L}\,
 \frac{\sinh(n\pi y/L)}{\sinh(n\pi H/L)},\qquad n\ge1.
\]
El cociente normaliza el valor en el lado que aporta el dato.
La continuidad hasta una esquina exige que los valores de sus lados
coincidan. Del mismo modo, un dato inicial incompatible con el contorno
puede impedir continuidad en la esquina inicial sin impedir una solución
para $t>0$ que recupere el dato interior o en norma cuadrática.
Consultas: soluciones de \exref{M2-E0109} y \exref{M2-E0268}.

%
\subsection*{Mallas pequeñas y alcance de las garantías}
\label{CH10-B04}
En calor, sean $x_i=i\Delta x$, $t_n=n\Delta t$ y
$r=\kappa\Delta t/(\Delta x)^2$. El esquema explícito es
\[
 U_i^{n+1}=rU_{i-1}^n+(1-2r)U_i^n+rU_{i+1}^n.
\]
Todos los valores de la derecha pertenecen al mismo nivel temporal.
Si $0\le r\le1/2$, es un promedio convexo: preserva no negatividad
con datos y contornos no negativos y controla perturbaciones en la norma
máxima con los mismos contornos. Esa propiedad no asegura precisión.
En una malla cuadrada de paso igual en ambas direcciones, Laplace da
$4U_{i,j}=U_{i-1,j}+U_{i+1,j}+U_{i,j-1}+U_{i,j+1}$.

Compárense nodos y tiempos coincidentes con una referencia. Si un modo
tiene multiplicador de un paso $G$, después de $N$ pasos su amplitud se
multiplica por $G^N$; no hacen falta $N$ actualizaciones manuales.
Al refinar, revísese la razón $r$. La coincidencia nodal no demuestra
exactitud continua, y observar un modo decreciente no prueba estabilidad
para todos los modos. Consultas: soluciones de \exref{M2-E0355},
\exref{M2-E0356} y \exref{M2-E0380}.

%
\subsection*{Transformar el tiempo en un problema de calor}
\label{CH10-B05}
Sea $U(x,s)=\mathcal L_t\{u(x,t)\}(s)$ y $f(x)=u(x,0)$.
Cuando se pueden intercambiar la transformación y las derivadas
espaciales, $u_t=\kappa u_{xx}$ se convierte en
\[
 \kappa U_{xx}-sU=-f(x).
\]
El dato inicial permanece en la EDO espacial. Transfórmense también
los contornos y resuélvase el problema completo antes de invertir.
Para contornos homogéneos de valor y $s>0$, la unicidad puede comprobarse
restando dos soluciones. Si $W$ es su diferencia, se multiplica
$\kappa W''-sW=0$ por $W$ y se integra por partes: los términos de borde desaparecen y
las integrales de los cuadrados obligan a que la diferencia sea cero.

Después de invertir, verifíquense ecuación, datos y convergencia de la
transformada. Esta conexión utiliza EDO de contorno y Laplace del
capítulo 6; no proporciona por sí sola una teoría para dominios no
acotados o para estabilidad de EDP. Consulta: solución de
\exref{M2-E0357}.
